\subsection{An actual two-dimensional construction excluding three classes of countable quotients}

\begin{claim}{A.28}{Proved}{Continuum-sized diagonalization within dimension two}
Let $B$ be a specified connected countable CW classifying
space of dimension at most two, with nontrivial perfect
fundamental group $H$. Put $\mathfrak c=2^{\aleph_0}$.
There is an actual CW inclusion $i:B\hookrightarrow X$
such that $X$ is aspherical, $\dim X\leq2$, and $X$ has
at most $\mathfrak c$ cells. It induces an embedding
$H\hookrightarrow G:=\pi_1X$ and an ordinary integral
homology isomorphism. The displayed copy of $H$ normally
generates $G$.

Every one-variable equation over $G$ whose total
variable exponent is $1$ or $-1$ has a solution
in $G$. Every quotient of $G$ has this same
equation-solving property.

The group $G$ has no nontrivial countable quotient in
any of the following classes:
\begin{enumerate}
\item mixed-identity-free groups;
\item torsion-free groups with exactly two conjugacy classes;
\item nonabelian uniformly simple groups, with a finite
uniform bound allowed to depend on the target.
\end{enumerate}
Thus $G$ is perfect, and an acyclic starting space $B$
gives an acyclic $X$. For a starting group with no
nontrivial finite images, $G$ also has no nontrivial
finite images.

In particular there is an acyclic map
\[
 t:BG\longrightarrow S^2,
 \qquad \dim BG=2,\qquad |G|\leq\mathfrak c,
 \tag{A.28.1}
\]
whose group has no nontrivial images in torsion groups
or in locally linear groups, and no
countable quotient in the three listed classes. The
initial two-dimensional resolution of $S^2$ may be
retained as an actual subcomplex, and $t$ may be chosen
to extend its specified resolution map literally.
No assertion that all countable quotients have been
excluded is made.
\end{claim}

\begin{proof}
Write $\mathcal C$ for the union of the three target
classes in the statement. We first verify the exact
subgroup version of the block attachment needed for
the cardinal bound.

\paragraph{A block using coefficients in a countable subgroup.}
Let $Z$ be an actual aspherical two-dimensional CW
complex, put $A=\pi_1Z$, and let $K\leq A$ be a
countable subgroup with an epimorphism $q:K\twoheadrightarrow E$
onto some $E\in\mathcal C$. Take the unsolvable
countable system over $E$ constructed in A.27. For
every occurring coefficient choose a lift in $K$ and
a loop in $Z$ representing that lift. Use the same
lift at each occurrence of a specified coefficient.
Attach its variable loops and two-cells over $Z$,
using these lifted words. Denote the resulting complex
by $Z'$.

This construction does not require a map $A\to E$.
Its geometric verification uses only the displayed
words over the torsion-free group $A$. In the chain
case (A.27.2), elimination of the non-root variables
in a finite prefix leaves one unimodular root relator
over the free product of $A$ with the tail free group.
In the tree case (A.27.8), elimination in descending
depth has the same outcome. The root exponent is
one because its defining outer word has exponent
zero in each substituted child. This holds for the
chosen lifts regardless of the relations between
their images in $E$. The Forester--Rourke theorem
and the finite-prefix argument in A.27 therefore
prove that $Z'$ is aspherical and that
$A\hookrightarrow\pi_1Z'$ is injective. The actual
relative ordinary boundary is the identity in
degrees two and one, so $Z\hookrightarrow Z'$ is
an integral homology isomorphism. It adds countably
many one- and two-cells.

The normal-generation computation is also unchanged.
Killing $A$ makes the root equation force its root
variable to be $1$; then every other equation forces
its own variable to be $1$. Thus $A$ normally
generates $\pi_1Z'$. Finally, any homomorphism
$\pi_1Z'\to E$ whose restriction to $K$ equals $q$
would evaluate the lifted words to the unsolvable
system over $E$. Consequently
\[
 \text{no homomorphism }\pi_1Z'\longrightarrow E
 \text{ restricts to }q\text{ on }K.
 \tag{A.28.2}
\]
This establishes the required block even when $q$
had no extension to $A$ in the first place.

\paragraph{One pass through all countable subgroup maps.}
Suppose $Z_\beta$ has at most $\mathfrak c$ cells
and put $A_\beta=\pi_1Z_\beta$. Choose one
representative of each group in $\mathcal C$ on
underlying set $\mathbb N$ or a finite initial
segment of $\mathbb N$. There are at most
$\mathfrak c$ possible group structures on these
sets. List every triple
\[
 (K,E,q),\qquad K\leq A_\beta\text{ countable},
 \quad E\in\mathcal C,
 \quad q:K\twoheadrightarrow E.
 \tag{A.28.3}
\]
There are at most
\[
 |A_\beta|^{\aleph_0}\leq
 \mathfrak c^{\aleph_0}=\mathfrak c
 \tag{A.28.4}
\]
countable subgroups, and at most $\mathfrak c$
functions between any two countable sets. Hence
(A.28.3) is a set of cardinality at most
$\mathfrak c$.

In this same pass also list every one-variable
word over $A_\beta$ with total variable exponent
$1$ or $-1$. There are at most $\mathfrak c$
such finite words. For each word adjoin one
fresh generator and its specified relator.
Forester--Rourke, in the form proved applicable
in A.18, gives an actual aspherical extension
of the current coefficient space and an
injective group map. The relative ordinary
boundary is multiplication by $1$ or $-1$,
so this is an integral homology isomorphism.
After killing the current coefficient group,
the relation is $x^{\pm1}=1$; hence the new
group is normally generated by that coefficient
group. These additional attachments are actual
one-cell and two-cell pairs and preserve all
the bounds needed for the pass.

Well-order the combined set of listed subgroup
maps and unimodular words, and attach the
corresponding blocks or one-pair extensions
successively.
The coefficients for a listed block are always the
chosen elements of its subgroup $K$ of the group
at the beginning of the pass. Every earlier
attachment embeds its coefficient group, so those
elements remain available in the current group
when this block is attached. The subgroup argument
above applies there; it does not assume that the
listed map $q$ extends over earlier attachments
within the pass.

At any intermediate limit take the actual CW union.
Every sphere map of dimension at least two has
compact image in a finite subcomplex and therefore
lies in an earlier aspherical stage. It is
nullhomotopic there, proving asphericity of the
union. Similarly, a nullhomotopy of an earlier
loop lies in one earlier stage, where the
coefficient injections preserve its nontriviality.
Thus the groups form a union of embedded subgroups.
Cellular chain complexes commute with these unions,
so homology isomorphisms persist at limits.
Normal generation persists as well: in any quotient
killing the group at the beginning of the pass,
the root and then all the other variables of
each successive block are forced to be trivial.

At most $\mathfrak c$ countable blocks have been
attached. Their union, called $Z_{\beta+1}$,
still has at most $\mathfrak c$ cells. It has
the same ordinary integral homology as $Z_\beta$,
and each listed map is obstructed by its actual
attached block through (A.28.2).

\paragraph{The outer union.}
Start with $Z_0=B$. Perform one such pass for each
$\beta<\omega_1$, taking actual unions at limit
ordinals, and put
\[
 X=\bigcup_{\beta<\omega_1}Z_\beta,
 \qquad G=\bigcup_{\beta<\omega_1}A_\beta.
 \tag{A.28.5}
\]
Since $\aleph_1\leq\mathfrak c$, this complex
has at most $\mathfrak c$ cells. The successor
and limit verifications give its asphericity,
dimension, homology, coefficient embedding, and
normal generation by $H$. In particular each
intermediate coefficient group is torsion-free,
as required when attaching the next block.

Suppose there were an epimorphism $p:G\twoheadrightarrow E$
with $E\in\mathcal C$. Choose a preimage in $G$
for each of the countably many elements of $E$.
Every preimage lies in an earlier $A_\beta$;
regularity of $\omega_1$ puts all these preimages
in a single $A_\gamma$, with $\gamma<\omega_1$.
The subgroup $K$ they generate is countable and
$p|_K$ is onto $E$. After identifying $E$ with
its chosen representative, the triple
$(K,E,p|_K)$ occurred in pass $\gamma$. Restricting
$p$ to the block attached for that triple
contradicts (A.28.2). This proves the asserted
exclusion of quotients.

To verify the added equation-solving assertion,
take a one-variable unimodular word over $G$.
Its finite coefficient list belongs to some
$A_\beta$. That exact word was included in
pass $\beta$, and its new generator is a
solution in $G$. If $p:G\twoheadrightarrow Q$
is any epimorphism and a unimodular word is
given over $Q$, choose lifts in $G$ of its
finitely many coefficients. The lifted word
has unchanged variable exponent. Map a
solution in $G$ through $p$ to obtain the
required solution in $Q$. This argument makes
no assertion about geometric attachments for
arbitrary finite nonsingular systems.

Perfectness and the acyclic-base assertion follow
from the homology isomorphism. If $H$ has no
nontrivial finite images, a map $G\to F$ into
a finite group kills $H$. Since $H$ normally
generates $G$, it kills all of $G$.

\paragraph{The specified map to the two-sphere.}
Take the finite two-dimensional classifying
resolution $t_0:B\to S^2$ from A.10. Its group
is perfect, its first homology is zero, its
second homology is $\mathbb Z$, and it has no
nontrivial finite images. Apply the construction
just proved to this actual $B$.

We verify the stronger image restrictions in the
statement directly. In the Higman presentation used
in A.10, suppose all four generator images have
finite order, in an otherwise arbitrary group.
If one image is nontrivial, choose the smallest
prime $p$ dividing one of those four orders, say
the order of $x$. Its preceding generator $y$
satisfies $y^{-1}xy=x^2$. Hence the order of $x$
is odd, and the multiplicative order of $2$ modulo
$p$ divides the finite order of $y$. That
multiplicative order is greater than one and
divides $p-1$, so a prime smaller than $p$ divides
the order of $y$, a contradiction. Thus every
homomorphism from Higman's group to a torsion
group is trivial; finiteness of the target was
not needed in this calculation.

Every finite-dimensional linear image of Higman's
group is trivial as well. Extend the coefficient
field to an algebraic closure. Since each generator
matrix $x$ is conjugate to $x^2$, squaring permutes
its finite multiset of nonzero eigenvalues. Every
eigenvalue is therefore a root of unity. In positive
characteristic a suitable power of $x$ is unipotent,
and a further power of the characteristic makes
it the identity. Thus all generator matrices have
finite order in this case.

In characteristic zero, choose positive integers
$L,m$ with $x^L$ and $y^m$ unipotent, where
$y^{-1}xy=x^2$ as above. The finite nilpotent
logarithm $v=\log(x^L)$ satisfies
\[
 \operatorname{Ad}(y^{-m})(v)=2^m v.
 \tag{A.28.7}
\]
Here the equality follows by applying the logarithm
polynomial to
$y^{-m}x^Ly^m=(x^L)^{2^m}$. Conjugation by the
unipotent matrix $y^{-m}$ is a unipotent linear
operator on the matrix algebra, so its only
eigenvalue is $1$. Since $2^m\ne1$, (A.28.7)
forces $v=0$; the inverse finite exponential
polynomial then gives $x^L=1$. Thus all generator
images again have finite order. The minimal-prime
argument proves that the entire linear image
is trivial.

Call a group locally linear here when every
finitely generated subgroup admits a faithful
finite-dimensional representation over some
field. An image of the four-generator Higman
group in such a target lies in one such subgroup,
so the preceding argument kills it. A.10 proves
that its group over the simply connected target
$S^2$ is normally generated by the inserted
copies of Higman's group: killing those copies
gives exactly the fundamental group of $S^2$.
Consequently the initial $H$ has no nontrivial
images in torsion or locally linear groups.
Since $H$ normally generates $G$, the same
conclusion holds for $G$.

Extend $t_0$ over the construction as follows.
Send every new loop to the basepoint of $S^2$.
The image of every two-cell attaching loop is
nullhomotopic because $S^2$ is simply connected;
choose a filling disk and use it to extend the
map over that cell. Do this at each successive
attachment and use the CW weak topology at
unions. This gives an actual continuous map
$t:X\to S^2$ with the strict equality
\[
 t i=t_0.
 \tag{A.28.6}
\]
Since $i$ and $t_0$ induce integral homology
isomorphisms, so does $t$. In particular its
second cohomology class is the extension of
the original generator: under
$H^2(X;\mathbb Z)\cong H^2(B;\mathbb Z)$ it
restricts to $t_0^*[S^2]$. The fundamental-group
map to $S^2$ is onto, and every local system
on $S^2$ is constant. The integral homology
isomorphism gives the isomorphism with every
constant abelian coefficient group by the
universal coefficient theorem. Thus $t$ is
acyclic in the precise sense used in A.2.
The nonzero second homology makes $\dim X=2$.
This proves (A.28.1).
\end{proof}

The quotient conclusion has an explicit scope.
For a homomorphism $G\to T$ to a countable group,
it proves triviality whenever its nontrivial image
would belong to one of the three listed classes.
No inheritance of these properties by arbitrary
subgroups of $T$ is used. In particular A.28
does not establish R2 for all countable target
groups at $W=S^2$.

There is a further exact restriction in the
$S^2$ example. Its initial group is finitely
generated, since A.10 supplies a finite model.
The resulting $G$ is therefore normally generated
by finitely many specified elements. Every
nontrivial quotient of $G$ has a maximal proper
normal subgroup: the union of a chain of proper
normal subgroups is proper, since a chain member
would contain all finitely many normal generators
if their images all lay in the union. Zorn's
lemma then supplies the maximal subgroup. Hence
every nontrivial countable quotient of this $G$
has an infinite countable simple quotient, since
the finite images of $G$ were already excluded.
Such a simple
quotient must satisfy a nontrivial mixed identity,
must fail every finite bound of uniform simplicity,
and cannot be a torsion-free two-conjugacy-class
group, by the proved exclusion above. It must also
contain an element of infinite order and fail
local linearity, by the direct Higman-image
calculations. It must solve every one-variable
unimodular equation over itself, by the
quotient permanence just proved.
